% English draft based on sections/11_instability_region_ja.tex
\section{The Instability Region: A Three-Axis Phase Diagram}

This chapter integrates the three main logical lines of the paper: non-closure and external stopping, threshold mobility and sapiens-like acceleration, and the three-layer dynamics $L_F/L_R/L_E$.

\subsection{Three Axes}

The dynamical state of intelligence is described by three axes:
\begin{itemize}
\item inferential velocity $v_{\text{inf}}$;
\item resolution $r_{\text{res}}$;
\item stopping hardness $\eta$.
\end{itemize}

Higher inferential velocity means more inferential steps per unit time. Higher resolution means finer distinctions and causal decompositions. Higher stopping hardness means less revisability of stopping.

In addition to these three axes, the transition into divergence depends on the log-reference velocity $v_{LR}$. This is the rate at which the log layer $L_R$ can recursively reference traces generated by $L_F$ and incorporate them as usable patterns. The difference
\[
\Delta v_{\text{intel}} = v_{\text{inf}} - v_{LR}
\]
measures the speed gap between inference generation and log reference. When this gap grows too large, inference outruns the log layer and the basis for $L_E$ begins to break.

\subsection{Three Regions}

The phase space contains three regions.

\textbf{Stable operation region.} Inference is constrained by externalized stopping, $L_E$ remains revisable, and $\mu_\theta$ can operate.

\textbf{Divergent region.} High velocity and high resolution combine with internalized stopping. $L_F$ runs ahead, $L_R$ cannot follow, $L_E$ cannot emerge, and $\mu_\theta$ cannot open a window.

\textbf{Rigid region.} Stopping hardness is excessive. $L_E$ fixes, $\mu_\theta$ loses mobility, and the dynamic relation between $L_F$ and $L_E$ freezes.

\begin{figure}[H]
\centering
\includegraphics[width=\linewidth]{fig8_phase_diagram.pdf}
\caption{The three-axis phase diagram of intelligence. Stable operation exists only in a narrow band between divergence and rigidity in the space of $v_{\text{inf}}$, $r_{\text{res}}$, and $\eta$.}
\label{fig:part2-phase-diagram}
\figurenote{In the stable region, $\mu_\theta$ functions and $L_E$ remains revisable. In the divergent region, $L_F$ runs away; in the rigid region, $L_E$ fixes.}{Instability is not a defect but a constitutive feature. If speed and resolution exceed the bandwidth of log reference and external stopping, intelligence falls into divergence.}
\end{figure}

\subsection{Four-Fold Collapse}

The divergent region is a four-fold collapse:
\begin{enumerate}
\item $L_F$ accelerates and generates too much inference;
\item $L_R$ cannot recursively reference the generated logs fast enough, because $\Delta v_{\text{intel}} = v_{\text{inf}} - v_{LR}$ has become critically large;
\item $L_E$ cannot emerge because its basis in $L_R$ is lost;
\item $\mu_\theta$ cannot open the reconfiguration window.
\end{enumerate}

This collapse is not caused by intelligence being too good. It is caused by high-speed, high-resolution inference being cut off from externalized stopping.

\subsection{Rigidity}

The opposite failure is rigidity. Here stopping exists, but it becomes too hard. A fixed $L_E$ prevents threshold mobility and blocks adaptation. Perfect closure would not be the completion of intelligence; it would be the end of intelligence.

\subsection{Instability as Constitutive}

The instability region does not mean the failure of intelligence. It is the boundary region that necessarily appears when finite intelligence with stopping structures attempts to maintain free inference. Religion, science, institutions, and communities can all be understood as stopping-operation systems for managing this instability region.

\bigskip
